\clearpage
\section{Problem Analysis}

The assignment is concurrency-oriented because the scan is naturally
asynchronous and because the implementation must expose three different
programming styles. The computation is simple in principle: visit every regular
file, read its size, increment one total counter, and increment one histogram
band. The difficulty comes from doing this over a real filesystem while keeping
the caller responsive.

\subsection{Filesystem and Concurrency Constraints}

The filesystem adds several sources of complexity:

\begin{itemize}[leftmargin=*]
  \item traversal is recursive and may reach deep directory hierarchies;
  \item the scan must remain responsive while data is collected;
  \item filesystem state may change while the scan is running;
  \item invalid paths, unreadable directories, symbolic links, and special files
  must not crash the program;
  \item the same logical result must be produced by all three paradigms.
\end{itemize}

The scan is mainly I/O-bound. Parallelism can improve responsiveness and can
overlap independent directory operations, but it also introduces coordination
costs. A correct implementation must therefore define when a directory is
considered visited, how partial counters are updated, how completion is
detected, and what cancellation means when work is already in flight.

\subsection{Correctness Properties}

The expected behaviour can be summarized through four properties:

\begin{itemize}[leftmargin=*]
  \item \textbf{Safety}: only regular files are counted and every counted file
  contributes to exactly one size band;
  \item \textbf{Consistency}: the three implementations must produce the same
  report on the same stable directory tree;
  \item \textbf{Responsiveness}: progress updates and GUI interactions must not
  be blocked by the traversal;
  \item \textbf{Clean termination}: completion, errors, and cancellation must
  release the waiting CLI or restore the GUI state.
\end{itemize}

The design keeps the report model independent from the execution strategy. This
makes the three implementations comparable while still allowing each engine to
use its own scheduling style.

\subsection{Shared Data Model}

The shared package \texttt{pcd.fsstat.common} contains the minimal common
model:

\begin{itemize}[leftmargin=*]
  \item \texttt{FSReport}: immutable report object containing directory, size
  threshold, band count, band histogram, total file count, and elapsed time;
  \item \texttt{FSReportListener}: callback contract for asynchronous updates,
  completion, and errors;
  \item \texttt{FSReportJob}: cancellation handle for imperative scanners;
  \item \texttt{SizeUnit}: user-facing unit conversion and formatting helper.
\end{itemize}

The report stores the band counts in a defensive copy so that callers cannot
mutate it accidentally.

\subsection{Band Semantics}

The range from zero to the maximum file size threshold is divided into a
configured number of regular bands, and an additional overflow band stores
files larger than that threshold. Boundary values are handled deterministically:
a file exactly equal to \texttt{maximumFileSizeBytes} belongs to the last
regular band, while \texttt{maximumFileSizeBytes + 1} belongs to the overflow
band. This rule is centralized in \texttt{FSReport.getBandIndex()} and is
tested explicitly with boundary-sized files.
